\subsection{Invariant-block quadratic rates and an exact small-gain criterion}
An invariant-block quadratic isolates a setting where differentiated scalar
weights can improve conditioning exactly. The result describes the benefit
available when the flag aligns with Hessian invariant subspaces. Learning a
flag from gradient covariance does not itself ensure this alignment; the
perturbation bound quantifies how flag error can erode the ideal rate.

\begin{assumption}[Invariant Hessian blocks]\label{ass:rigid-quadratic}
\begin{gather}
 F(w)=\tfrac12(w-w^\star)^\top H(w-w^\star),\quad H\succ0,\\
 H\Pi_j=\Pi_jH,\quad
 0<\ell_j\Pi_j\preceq\Pi_jH\Pi_j\preceq u_j\Pi_j,\quad
 B=\sum_{j=1}^{K+1}a_j\Pi_j,\quad a_j>0.
\end{gather}
The projectors in this assumption act on the parameter vector; a matrix
right-action uses $I_m\otimes\Pi_j$ after vectorization.
\end{assumption}
\begin{theorem}[Block conditioning and optimal scalar weights]\label{thm:rigid-conditioning}
Under Assumption~\ref{ass:rigid-quadratic}, let
\begin{gather}
 \ell_B:=\min_j a_j\ell_j,\quad u_B:=\max_j a_ju_j,\quad
 \kappa_B:=u_B/\ell_B,\quad \eta_\star:=2/(u_B+\ell_B),\\
 w_{t+1}=w_t-\eta B\nabla F(w_t).
\end{gather}
Then
\begin{align}
 \|w_t-w^\star\|_{B^{-1}}
 &\le\left[\max_{\lambda\in[\ell_B,u_B]}|1-\eta\lambda|\right]^t
                  \|w_0-w^\star\|_{B^{-1}},\\
 F(w_t)&\le\left(\frac{\kappa_B-1}{\kappa_B+1}\right)^{2t}F(w_0)
                 \quad(\eta=\eta_\star),\\
 \inf_{a_1,\ldots,a_{K+1}>0}\frac{\max_j a_ju_j}{\min_j a_j\ell_j}
 &=\max_j\frac{u_j}{\ell_j},\qquad a_j=1/\ell_j\text{ attains the infimum}.
\end{align}
If $H=\sum_jh_j\Pi_j$ and $a_j=h_j^{-1}$, then $w_1=w^\star$ for $\eta=1$.
\end{theorem}
\begin{proof}
\begin{align}
 x_t&:=B^{-1/2}(w_t-w^\star),\quad
 \widetilde H:=B^{1/2}HB^{1/2}=BH,\\
 x_{t+1}&=(I-\eta\widetilde H)x_t,\quad
 \operatorname{spec}(\widetilde H)\subset[\ell_B,u_B],\\
 F(w_t)&=\tfrac12x_t^\top\widetilde Hx_t,\quad
 \widetilde H^{1/2}x_t=(I-\eta\widetilde H)^t\widetilde H^{1/2}x_0,\\
 \min_{\eta>0}\max\{|1-\eta\ell_B|,|1-\eta u_B|\}
 &=\frac{u_B-\ell_B}{u_B+\ell_B},\\
 \frac{\max_i a_iu_i}{\min_i a_i\ell_i}
 &\ge\frac{a_ju_j}{a_j\ell_j}=u_j/\ell_j\quad(\forall j),\\
 a_j=1/\ell_j&\Longrightarrow\min_j a_j\ell_j=1,\quad
                                  \max_j a_ju_j=\max_j u_j/\ell_j.
\end{align}
\end{proof}
\begin{proposition}[Flag perturbation of quadratic conditioning]\label{prop:rigid-quadratic-error}
Let $B_\star=B(U^\star)$, $B=B(U)$, and
$\delta_B:=\Gamma^a d_P(U,U^\star)$. If
$\ell_\star I\preceq H^{1/2}B_\star H^{1/2}\preceq u_\star I$, then
\begin{gather}
 \ell_\star-\|H\|_{\op}\delta_B
 \le\lambda_{\min}(H^{1/2}BH^{1/2})
 \le\lambda_{\max}(H^{1/2}BH^{1/2})
 \le u_\star+\|H\|_{\op}\delta_B,\\
 \|H\|_{\op}\delta_B<\ell_\star
 \Longrightarrow\quad
 \kappa(H^{1/2}BH^{1/2})
 \le\frac{u_\star+\|H\|_{\op}\delta_B}
            {\ell_\star-\|H\|_{\op}\delta_B}.
\end{gather}
\end{proposition}
\begin{proof}
\begin{equation}
 \|H^{1/2}(B-B_\star)H^{1/2}\|_{\op}
 \le\|H\|_{\op}\|B-B_\star\|_{\op}\le\|H\|_{\op}\delta_B.
\end{equation}
\end{proof}
The small-gain criterion addresses a different question: when does direct
parameter contraction dominate feedback through optimizer state? The resolvent
$(I-D)^{-1}$ sums repeated propagation inside the state subsystem. Its coupling
cost must remain below the direct contraction margin $\alpha$, giving an exact
criterion for the stated nonnegative envelope.

\begin{theorem}[Schur-resolvent small gain]\label{thm:rigid-schur}
This is the nonsingular $M$-matrix/Schur-complement criterion for a nonnegative
feedback system~\citep{berman1994nonnegative}, specialized to the optimizer's
parameter--state partition.
Let a nonnegative sensitivity envelope at a fixed step size have the form
\begin{gather}
 \mathcal A_\eta:=
 \begin{pmatrix}1-\eta\alpha&\eta b^\top\\c&D\end{pmatrix}\ge0,\quad
 b,c\ge0,\quad D\ge0,\quad\rho(D)<1,\quad\eta>0.
\end{gather}
Then
\begin{equation}
 \rho(\mathcal A_\eta)<1\quad\Longleftrightarrow\quad
                      \alpha>b^\top(I-D)^{-1}c.
\end{equation}
If the right side holds, choose $\varepsilon>0$ so small that
\begin{gather}
 v^\top:=(b^\top+\varepsilon\mathbf1^\top)(I-D)^{-1}>0,\quad
 \delta:=\alpha-v^\top c>0,\quad a_\eta:=(1,\eta v^\top),\\
 0<\gamma\le\delta,\quad
 \gamma\eta\max_jv_j\le\varepsilon.
\end{gather}
Then
\begin{equation}
 a_\eta\mathcal A_\eta\preceq(1-\gamma\eta)a_\eta.
\end{equation}
\end{theorem}
\begin{proof}
\begin{align}
 (I-D)^{-1}&=\sum_{k=0}^\infty D^k\ge0,\\
 v^\top(I-D)&=b^\top+\varepsilon\mathbf1^\top,\\
 a_\eta\mathcal A_\eta
 &=(1-\eta(\alpha-v^\top c),\ \eta v^\top-\eta\varepsilon\mathbf1^\top)
 \preceq(1-\gamma\eta)a_\eta.
\end{align}
Thus $\rho(\mathcal A_\eta)\le1-\gamma\eta<1$. Conversely,
$\rho(\mathcal A_\eta)<1$ implies the existence of $x=(x_0,x_s)>0$ with
$(I-\mathcal A_\eta)x>0$, e.g.
$x=(I-\mathcal A_\eta)^{-1}\mathbf1$. Hence
\begin{gather}
 (I-D)x_s>cx_0,\quad x_s>(I-D)^{-1}cx_0,\quad
 \alpha x_0>b^\top x_s\ge b^\top(I-D)^{-1}cx_0.
\end{gather}
\end{proof}
\begin{corollary}[A uniform $1/n$ sensitivity floor]\label{cor:rigid-schur-stability}
Suppose a coupled optimizer satisfies
\begin{gather}
 \binom{u_{t+1}}{s_{t+1}}\preceq
 \mathcal A_\eta\binom{u_t}{s_t}
 +\frac{I_t}{b_t}\binom{\eta d}{h},\quad
 \PP(I_t=1)=b_t/n,\quad (u_0,s_0)=0,
\end{gather}
with the hypotheses of Theorem~\ref{thm:rigid-schur}. Then
\begin{align}
 \EE[u_T+\eta v^\top s_T]
 &\le\frac{d+v^\top h}{n\gamma}
                         (1-(1-\gamma\eta)^T),\\
 \epsilon_T^{\mathrm{stab}}&\le\frac{L_0(d+v^\top h)}{n\gamma}.
\end{align}
\end{corollary}
\begin{proof}
\begin{gather}
 \EE[u_{t+1}+\eta v^\top s_{t+1}]
 \le(1-\gamma\eta)\EE[u_t+\eta v^\top s_t]
                                      +\eta(d+v^\top h)/n.
\end{gather}
\end{proof}
\begin{proposition}[Local Jacobian envelopes for the actual map]\label{prop:rigid-jacobian}
Let $\mathcal T_{t,S}$ denote the complete map of
Definition~\ref{def:rigid-map}, expressed in regular local coordinates. For a convex
coupling region $\mathcal K$, partition its Jacobian into parameter and state blocks:
\begin{gather}
 \mathrm D\mathcal T_{t,S}(x)=
 \begin{pmatrix}J_{ww}&J_{ws}\\J_{sw}&J_{ss}\end{pmatrix},\quad
 \sup_{x\in\mathcal K}\|J_{ww}\|\le1-\eta\alpha,\quad
 \sup_{x\in\mathcal K}\|J_{ws}\|\le\eta b,\\
 \sup_{x\in\mathcal K}\|J_{sw}\|\le c,\quad
 \sup_{x\in\mathcal K}\|J_{ss}\|\le q<1.
\end{gather}
If $\alpha(1-q)>bc$, the scalar-state version of
Theorem~\ref{thm:rigid-schur} applies. Block norms give the matrix-state version.
For the broad triangle-inequality envelope of
Definition~\ref{def:rigid-transition}, these contraction conditions are additional
hypotheses, not consequences of smoothness alone.
\end{proposition}
\begin{proof}
\begin{equation}
 \mathcal T_{t,S}(x)-\mathcal T_{t,S}(x')
 =\int_0^1\mathrm D\mathcal T_{t,S}(x'+a(x-x'))(x-x')\,\mathrm da.
\end{equation}
\end{proof}
